% TOP_PROBLEM: {"id": 3100061, "title": "Is x^(2)+y^(2)=z^(2) partition regular", "category_id": 1, "category": {"id": 1, "name": "number_theory", "display_name": "Number Theory", "description": "Properties of integers, prime numbers, Diophantine equations.", "slug": "number-theory", "order_index": 1, "created_at": "2026-07-31T15:26:25.670Z"}, "status": "partially_solved", "problem_number": "AMR-030-0061", "set_id": 13, "statusQualification": "The established two-color case is separated from partition regularity for an arbitrary finite number of colors; no primitivity condition is imposed."}
% FIRST_POSTED: 2026-10-02
% ENTRY_KIND: statement_only
% UnsolvedMath ID 3100061; source catalogue code AMR-030-0061
% !TeX program = lualatex
% Definition and sources only; this file is not an LLM solution attempt.
\documentclass[11pt,a4paper]{article}
\usepackage[margin=25mm]{geometry}
\usepackage{fontspec}
\setmainfont{lmroman10-regular.otf}[BoldFont=lmroman10-bold.otf,
 ItalicFont=lmroman10-italic.otf,BoldItalicFont=lmroman10-bolditalic.otf]
\usepackage{amsmath,amssymb,mathtools}
\usepackage{unicode-math}
\setmathfont{latinmodern-math.otf}
\AtBeginDocument{\let\setminus\smallsetminus}
\usepackage{xurl,hyperref}
\hypersetup{colorlinks=true,urlcolor=blue}
\setcounter{secnumdepth}{0}
\setlength{\parindent}{0pt}
\setlength{\parskip}{5pt}
\setlength{\emergencystretch}{3em}
\begin{document}
\section*{Is x\textasciicircum{}(2)+y\textasciicircum{}(2)=z\textasciicircum{}(2) partition regular}\label{3100061}
{\small UnsolvedMath ID 3100061 · AMR-030-0061 · Number Theory · AMR Open Problem Lists}
\subsection{Definitions and mathematical statement}
A finite coloring of the positive integers is a map
$c:\mathbb Z_{>0}\to\{1,\ldots,r\}$, where $r\ge1$ is finite
but otherwise arbitrary. A monochromatic Pythagorean triple is
a triple of positive integers $x,y,z$ satisfying
\[
                  x^2+y^2=z^2,
                    \qquad c(x)=c(y)=c(z).
\]
Is the Pythagorean equation partition regular over the positive
integers, meaning that every such finite coloring has a
monochromatic solution?

The three integers are necessarily distinct: positivity gives
$z>x,y$, and $x=y$ would make $z/x=\sqrt2$ rational. They need
not form a primitive triple; imposing $\gcd(x,y,z)=1$ would be
an additional restriction. The number of colors is fixed within
each coloring, but the assertion must hold for every finite $r$.

The two-color case has already been established and is not the
remaining general question. Likewise, conclusions making only
two of the three entries monochromatic do not establish full
partition regularity. Equivalently, for each $r$ one asks for a
finite $N(r)$ such that every $r$-coloring of $[N(r)]$ contains
the required triple; the finite and infinite formulations are
related by the usual compactness argument for finite colorings.
\subsection{Short English statement}
Does every finite coloring of the positive integers contain a monochromatic Pythagorean triple?
\subsection{Sources}
\begin{itemize}
\item[S1] ULAM.AI and the UnsolvedMath Contributors, supplied problems.json, record 3100061 (AMR-030-0061), AMR Open Problem Lists. Catalogue identity and source transcription; CC BY 4.0.
\url{https://huggingface.co/datasets/ulamai/UnsolvedMath}.
\item[S2] Cooper - Combinatorial Problems I Like (2020 snapshot), Combinatorial Number Theory, source-order bullet 61. Original source indexed by AMR-030-0061.
\url{https://people.math.sc.edu/cooper/combprob.html}.
\item[S3] Heule, Kullmann and Marek, Solving and Verifying the Boolean Pythagorean Triples problem.
\url{https://arxiv.org/abs/1605.00723}.
\end{itemize}
\end{document}
